\sgaStableTarget{sga6.e14.note-for-page-1}\subsection*{Note for page 1.}

The theory of the operations @@M0@@ in @@M1@@, and of the operations @@M2@@ in @@M3@@, was written down by P. Deligne in 1968. He proves that @@M4@@ thereby becomes a pre-@@M5@@-ring, and, when @@M6@@ is a sheaf of Algebras over a field, that @@M7@@ is even a @@M8@@-ring. It is plausible that a refinement of Deligne's method would make it possible to dispense with the restrictive hypothesis on @@M9@@.

\sgaStableTarget{sga6.e14.notes-for-page-11}\subsection*{Notes for page 11}

\noindent(*) (April 1969) D. MUMFORD points out to us a proof of (4.6) in the \uline{Chow ring} (without disregarding torsion), which he obtained in a seminar at Harvard in 1959! The reader will find the proof of this beautiful result in SGA 5 VII 9. This result obviously contains the corresponding cohomological result mentioned in note (**). However, the cohomological result is proved by J. P. JOUANOLOU under more general conditions (without a quasi-projectivity hypothesis, and in the relative smooth case over an arbitrary base).

\medskip

\noindent(**) (January 1969) J. P. JOUANOLOU has just proved these formulas for every @@M10@@ without disregarding torsion, cf. SGA 5 VII \S{} 4.9. DELIGNE and MANIN 17 had previously observed independently, and by (easy) arguments of a rather different nature, that it was enough to establish formulas for ``algebraic'' @@M11@@, and even for @@M12@@ in 16, to conclude the same relation in the general case.

% Final-index helper macros
\providecommand{\Qcoh}{\operatorname{Qcoh}}
\providecommand{\Coh}{\operatorname{Coh}}
\providecommand{\Parf}{\operatorname{Parf}}
\providecommand{\Lib}{\operatorname{Lib}}
\providecommand{\Loclib}{\operatorname{Loclib}}
\providecommand{\Proj}{\operatorname{Proj}}
\providecommand{\Mod}{\operatorname{Mod}}
\renewcommand{\Ob}{\operatorname{Ob}}
\providecommand{\Hom}{\operatorname{Hom}}
\providecommand{\toramp}{\operatorname{tor.amp}}
\providecommand{\parfamp}{\operatorname{par.amp}}
\providecommand{\coh}{\operatorname{coh}}
\providecommand{\qcoh}{\operatorname{qcoh}}
\providecommand{\parf}{\operatorname{parf}}
\providecommand{\torf}{\operatorname{torf}}
\providecommand{\amp}{\operatorname{amp}}
\providecommand{\num}{\operatorname{num}}
\providecommand{\naif}{\operatorname{naif}}
\providecommand{\alg}{\operatorname{alg}}
\providecommand{\Pic}{\operatorname{Pic}}
\providecommand{\Gr}{\operatorname{Gr}}
\providecommand{\Fil}{\operatorname{Fil}}
\providecommand{\Ch}{\operatorname{Ch}}
\providecommand{\Todd}{\operatorname{Todd}}
\providecommand{\Diff}{\operatorname{Diff}}
\providecommand{\Dcat}{\mathrm D}
\providecommand{\Tr}{\operatorname{Tr}}
\providecommand{\Z}{\mathbb Z}
\providecommand{\Q}{\mathbb Q}

\newpage
